\section{BCI 简史与神经科学基础}

\subsection{从大脑电信号的发现说起}

\begin{figure}[H]
\centering
\includegraphics[width=0.95\textwidth]{slides-images/slide-007.jpg}
\caption{Richard Caton 于 1875 年在《英国医学杂志》上发表的关于大脑电流的论文\protect\footnotemark}
\end{figure}
\footnotetext{Slides 第8页。这是最早关于脑电活动的科学记录之一。}

BCI 的历史可以追溯到 19 世纪。1875 年，英国科学家 \textbf{Richard Caton} 在动物实验中发现了大脑可以产生电信号，并且这些电信号与动物的行为（如转头、咀嚼）存在对应关系。这是人类首次通过实验证明``大脑以电信号编码信息''。

\subsection{脑电图（EEG）的发明}

\begin{figure}[H]
\centering
\includegraphics[width=0.95\textwidth]{slides-images/slide-008.jpg}
\caption{Hans Berger 发明 EEG 并于 1924 年记录了人类第一份脑电图\protect\footnotemark}
\end{figure}
\footnotetext{Slides 第9页。}

1924 年，德国精神科医生 \textbf{Hans Berger} 发明了\textbf{脑电图（Electroencephalogram, EEG）}——一种将电极放置在头皮表面来测量大脑电活动的技术。Berger 发现，大脑产生的电信号呈现出不同频率的波形，且这些波形与患者的精神状态密切相关：

\begin{itemize}
    \item \textbf{Alpha 波}（$\sim$8--13 Hz）：出现在安静闭眼时，代表放松状态
    \item \textbf{Beta 波}（$\sim$13--30 Hz）：出现在睁眼、进行认知任务时，代表活跃状态
\end{itemize}

\begin{knowledgebox}{Hans Berger 的故事}
Berger 原本是一名士兵。有一天他在训练中从马上摔落，遭受了脑震荡。而就在同一天，他的双胞胎姐妹突然感到不安，发电报向父亲询问弟弟的安危。这件事深深吸引了 Berger，促使他开始研究是否存在``心灵感应''——一种通过脑电波连接两个人的可能性。虽然心灵感应并未被证实，但他因此发明了 EEG，至今仍被广泛用于癫痫等疾病的诊断。
\end{knowledgebox}

\subsection{从头皮到大脑内部：侵入式记录}

EEG 虽然安全无创，但其空间分辨率非常有限。头皮上的电极测量的是\textbf{数百万神经元活动的平均值}——就像站在隔壁房间外试图听清里面的谈话，我们只能分辨出大致的情绪和话题，却无法听清具体内容。

\begin{figure}[H]
\centering
\includegraphics[width=0.95\textwidth]{slides-images/slide-010.jpg}
\caption{从头皮到单个神经元：将电极植入运动皮层（motor cortex）以获取高质量信号\protect\footnotemark}
\end{figure}
\footnotetext{Slides 第11页。左侧为 Hubel 1988 年的单细胞电极记录示意，右侧为运动皮层在大脑中的位置。}

为了获取更精确的信号，研究者们开始将电极\textbf{直接植入大脑内部}，紧邻单个神经元进行测量。本讲的研究主要关注\textbf{运动皮层（motor cortex）}——位于大脑中央沟前方的区域，负责控制全身肌肉的运动。

\begin{warningbox}{侵入式 vs 非侵入式记录的权衡}
侵入式记录（如微电极阵列）能获得单个神经元级别的高质量信号，但需要开颅手术植入电极，存在感染、组织损伤等风险。非侵入式记录（如 EEG、fMRI）安全无创，但信号质量低、分辨率差。当前临床 BCI 研究主要使用侵入式方法，因为高质量信号是实现精确解码的前提。
\end{warningbox}

\subsection{神经元的基本通信方式}

\begin{figure}[H]
\centering
\includegraphics[width=0.95\textwidth]{slides-images/slide-011.jpg}
\caption{神经元通过突触传递动作电位（spike）进行通信\protect\footnotemark}
\end{figure}
\footnotetext{Slides 第12页。来自 Goodman, Spiking Neural Networks。}

神经元是大脑的基本计算单元。每个神经元由\textbf{胞体（soma）}、\textbf{轴突（axon）}和\textbf{树突（dendrite）}组成。神经元之间通过\textbf{突触（synapse）}进行通信：当一个神经元需要向下游传递信息时，它会沿轴突产生一个\textbf{动作电位（action potential）}，也称为 \textbf{spike}——一个快速的电压脉冲。

如果将微电极放置在神经元轴突附近，就能记录到一系列尖锐的电压尖峰，即\textbf{脉冲序列（spike train）}。这就是 BCI 系统的原始输入数据。

\begin{importantbox}{神经元编码的核心特性}
\begin{itemize}
    \item 单个神经元的 spike 出现在毫秒级时间尺度上
    \item 同一神经元在相同实验条件下的 spike 模式具有\textbf{试次间变异性}（trial-to-trial variability）——这是生物神经元的固有噪声
    \item 信息不是由单个 spike 携带，而是由 spike 的\textbf{发放率（firing rate）}和\textbf{群体活动模式}共同编码
\end{itemize}
\end{importantbox}

\subsection{运动方向的神经编码}

\begin{figure}[H]
\centering
\includegraphics[width=0.95\textwidth]{slides-images/slide-012.jpg}
\caption{单神经元编码运动方向：猴子向左 vs 向右移动手臂时同一神经元的 spike 模式\protect\footnotemark}
\end{figure}
\footnotetext{Slides 第13页。来自 Shenoy \& Yu, Brain Machine Interfaces。}

经典实验表明，运动皮层中的单个神经元对运动方向具有\textbf{方向选择性（directional tuning）}。例如，某个神经元可能在猴子向右移动手臂时大量发放 spike（高发放率），而在向左移动时发放较少。

科学家发现，单个神经元的发放率与运动方向之间的关系可以用\textbf{余弦调谐曲线（cosine tuning curve）}来描述：

$$f(\theta) = b_0 + b_1 \cos(\theta - \theta_{\text{pref}})$$

其中：
\begin{itemize}
    \item $f(\theta)$：神经元在运动方向为 $\theta$ 时的发放率
    \item $\theta_{\text{pref}}$：该神经元的\textbf{偏好方向}（firing rate 最高的方向）
    \item $b_0$：基线发放率
    \item $b_1$：调制幅度
\end{itemize}

\subsection{多神经元群体解码}

\begin{figure}[H]
\centering
\includegraphics[width=0.95\textwidth]{slides-images/slide-013.jpg}
\caption{利用机器学习分类器从多神经元群体活动中解码运动方向\protect\footnotemark}
\end{figure}
\footnotetext{Slides 第14页。来自 Shenoy \& Yu, Brain Machine Interfaces。每种颜色代表一个运动方向，虚线为分类决策边界。}

由于单个神经元的信号具有固有噪声，仅靠一两个神经元无法精确判断运动方向。但如果同时记录\textbf{数百个神经元}的活动，利用机器学习分类器（如线性分类器、SVM 等）进行群体解码，就可以大幅提高解码精度。

这就是 BCI 的基本范式：
\begin{enumerate}
    \item 在运动皮层植入多电极阵列（Multi-Electrode Array, MEA）
    \item 记录数百个神经元的同步活动
    \item 训练机器学习模型，将神经活动模式映射到运动意图
    \item 在实时场景中使用训练好的解码器，将患者的``想法''转化为设备控制命令
\end{enumerate}

\subsection{神经活动测量技术概览}

\begin{figure}[H]
\centering
\includegraphics[width=0.95\textwidth]{slides-images/slide-014.jpg}
\caption{不同神经活动测量技术的时空分辨率对比\protect\footnotemark}
\end{figure}
\footnotetext{Slides 第15页。来自 Belliveau (1992)。x 轴为时间分辨率，y 轴为空间分辨率。}

不同的神经记录技术在\textbf{空间分辨率}和\textbf{时间分辨率}两个维度上存在巨大差异：

\begin{table}[H]
\centering
\begin{tabular}{lccc}
\toprule
\textbf{技术} & \textbf{空间分辨率} & \textbf{时间分辨率} & \textbf{是否侵入} \\
\midrule
fMRI & mm 级（脑区） & $\sim$0.5--1 秒 & 否 \\
EEG / MEG & cm 级（多脑区） & $\sim$毫秒 & 否 \\
ECoG & mm 级（皮层表面） & $\sim$毫秒 & 是（皮层表面） \\
微电极阵列 & $\mu$m 级（单神经元） & $\sim$毫秒 & 是（穿入皮层） \\
\bottomrule
\end{tabular}
\caption{各种神经记录技术特性对比}
\end{table}

理想情况下，BCI 需要\textbf{同时具备高空间分辨率和高时间分辨率}的记录技术。目前临床 BCI 研究中使用最多的是\textbf{微电极阵列}——一个约指甲大小的硅基芯片，上面排列着 96 根微型针状电极，每根电极可以测量附近几个神经元的活动，总共可以同时记录数百个神经元。

\subsection{本章小结}

\begin{itemize}
    \item BCI 技术的发展建立在 150 多年的神经科学研究基础上
    \item 运动皮层中的神经元以发放率编码运动方向信息，具有余弦调谐特性
    \item 单个神经元信号噪声大，需要多神经元群体解码来提高精度
    \item 微电极阵列是当前 BCI 临床研究的主要记录技术，提供单神经元级别的高时空分辨率信号
\end{itemize}

%% ========== Part 3: 从运动 BCI 到语音 BCI ==========

